\section{Conclusiones}

\begin{enumerate}
\item La corrección de densidad por altitud ($k$ = 0.733) sigue siendo el factor que gobierna la selección del ventilador en Cajicá: reduce la presión disponible y obliga a seleccionar el equipo en el punto equivalente de catálogo 3 840 m$^3$/h a 225 Pa (165 Pa en el sitio, escenario de filtro cargado).
\item El cambio de alcance del cliente elimina la presurización del cuarto. En consecuencia, el ventilador pasa de centrífugo a axial tubular: el punto de catálogo bajó de 260 Pa a 225 Pa, rango plenamente cubierto por axiales tubulares en PRFV, y se descartan el damper de alivio barométrico, el transmisor de presión diferencial y el manómetro local.
\item El punto de trabajo del sistema queda congelado en 165 Pa totales en el sitio para el escenario de filtro cargado (154 Pa de filtro + 11 Pa de pérdida de descarga en las rejillas), con potencia teórica de aire de 0.320 kW (0.43 HP) a $\eta$ = 0.55; el motor provisional es de 0.75 HP TEFC anticorrosivo, 440 V, 3$\phi$, 60 Hz, con margen de servicio 1.5 sobre el eje. La potencia definitiva se confirma con la curva del axial seleccionado.
\item La etapa definitiva de filtración es MERV 13-14 (prefiltro MERV 8 + filtro final V-bank de marco plástico); el laboratorio es de análisis industrial y no de biocontención, por lo que la filtración HEPA quedó expresamente descartada.
\item El ambiente exterior de la planta de hipoclorito de calcio exige PRFV viniléster o PP en la máquina de impulsión, inox 316L en rejillas y mallas, ferretería inox 316, sellante de silicona RTV neutra y ejecución anticorrosiva de la cubierta intemperie (PRFV o galvanizado G90 con pintura epóxica); el acero galvanizado desnudo queda descartado y el acero epóxico se limita a segunda línea con inspección programada.
\item La selección recomendada es un ventilador axial mural (placa mural) Ø560 mm en PRFV o ejecución epóxica anticorrosiva con transmisión directa, uniforme con el montaje típico instalado en la planta (Sodeca HQD/HGT mural anticorrosivo como primera opción de referencia; Greenheck mural, Aerovent/Twin City mural FRP, New York Blower FRP mural como alternativas), con el motor encapsulado severe duty dentro de la corriente de aire, el banco de filtración alojado en la cubierta intemperie, estructura de unión pernada al muro y malla de protección interior. Quedan pendientes de confirmación comercial el modelo y RPM exactos del ventilador con el software del fabricante; la tensión de la red (440 V, 3$\phi$, 60 Hz) y el plazo máximo de entrega ($\sim$3 meses) fueron confirmados por el cliente el 23/07/2026.
\end{enumerate}
